
\documentclass{article}
\usepackage{colortbl}
\usepackage{makecell}
\usepackage{multirow}
\usepackage{supertabular}

\begin{document}

\newcounter{utterance}

\centering \large Interaction Transcript for game `cladder', experiment `full\_v1.5\_default', episode 7517 with reflexion{-}Qwen3.5{-}9B.
\vspace{24pt}

{ \footnotesize  \setcounter{utterance}{1}
\setlength{\tabcolsep}{0pt}
\begin{supertabular}{c@{$\;$}|p{.15\linewidth}@{}p{.15\linewidth}p{.15\linewidth}p{.15\linewidth}p{.15\linewidth}p{.15\linewidth}}
    \# & \multicolumn{2}{c}{Player} && \multicolumn{2}{c}{Game Master} \\
    \hline

    \theutterance \stepcounter{utterance}  
    & & & \multicolumn{4}{p{0.6\linewidth}}{
        \cellcolor[rgb]{0.9,0.9,0.9}{
            \makecell[{{p{\linewidth}}}]{
                \texttt{\tiny{[P1$\langle$GM]}}
                \texttt{You have attempted this same task before and did not succeed. These are the plans you wrote after those attempts.} \\
\texttt{Trial \#1: New plan: I failed because I incorrectly identified the causal structure, mislabeling \$Vubr\$ as a confounder and \$Wibl\$ as a common cause of \$Zuph\$ and \$Vubr\$ when the prompt explicitly states \$Wibl\$ only affects \$vubr\$ and \$zuph\$ separately (no arrow between them). Consequently, I wrongly concluded that conditioning on \$Vubr\$ was necessary to block a backdoor path. In reality, since \$Wibl\$ is not a common cause of \$Zuph\$ and \$Vubr\$, there is no backdoor path between them; the only path is the direct causal chain \$Zuph \textbackslash{}rightarrow Uvzi\$. Therefore, Method 2 (unconditional correlation) correctly captures the direct effect, making Method 1 (conditioning on \$Vubr\$) incorrect as it would block the causal path. Next time, I will strictly map the arrows from the prompt first: \$Wibl \textbackslash{}to Zuph\$, \$Wibl \textbackslash{}to Vubr\$, \$Zuph \textbackslash{}to Uvzi\$, \$Vubr \textbackslash{}to Uvzi\$. I will then identify that \$Zuph\$ and \$Vubr\$ are independent given no common cause, meaning no adjustment is needed, and answer "No" with reasoning based on the absence of confounding.} \\
\\ 
\\ 
\texttt{You are an expert in causal inference. The following question is not a typical commonsense query, but rather a meticulously designed question created by a professor specializing in causal inference, intended to assess the students' mastery of the course content.} \\
\texttt{Imagine a self{-}contained, hypothetical world with only the following conditions, and without any unmentioned factors or causal relationships: Wibl has a direct effect on vubr and zuph. Zuph has a direct effect on uvzi. Vubr has a direct effect on uvzi. Method 1: We look at how zuph correlates with uvzi case by case according to vubr. Method 2: We look directly at how zuph correlates with uvzi in general. To understand how zuph affects uvzi, is it more correct to use the Method 1 than Method 2?} \\
\texttt{Start your answer with 'Yes' or 'No', followed by additional reasoning or evidence to support your explanation.} \\
            }
        }
    }
     \\ \\

    \theutterance \stepcounter{utterance}  
    & \multicolumn{4}{p{0.6\linewidth}}{
        \cellcolor[rgb]{0.561,0.737,0.561}{
            \makecell[{{p{\linewidth}}}]{
                \texttt{\tiny{[P1$\rangle$GM]}}
                \texttt{No} \\
\\ 
\texttt{The causal structure described is a directed acyclic graph (DAG) with the following arrows: \$Wibl \textbackslash{}to Zuph\$, \$Wibl \textbackslash{}to Vubr\$, \$Zuph \textbackslash{}to Uvzi\$, and \$Vubr \textbackslash{}to Uvzi\$. In this structure, \$Zuph\$ and \$Vubr\$ are not confounders of each other because there is no path connecting them that starts and ends with an arrow pointing away from the common cause \$Wibl\$ in a way that creates a backdoor path between \$Zuph\$ and \$Uvzi\$ via \$Vubr\$. Specifically, the path \$Zuph \textbackslash{}leftarrow Wibl \textbackslash{}to Vubr \textbackslash{}to Uvzi\$ is a fork at \$Wibl\$, but since \$Wibl\$ is an ancestor of both, conditioning on \$Vubr\$ does not open a backdoor path between \$Zuph\$ and \$Uvzi\$; rather, \$Vubr\$ is a mediator (or a parallel cause) on the path from \$Zuph\$ to \$Uvzi\$ only if we consider the total effect, but here \$Vubr\$ is a distinct cause of \$Uvzi\$ alongside \$Zuph\$.} \\
\\ 
\texttt{Crucially, to estimate the causal effect of \$Zuph\$ on \$Uvzi\$, we must consider the total effect, which includes the direct path \$Zuph \textbackslash{}to Uvzi\$. Method 1 conditions on \$Vubr\$ (looking at the correlation case by case according to \$Vubr\$). Since \$Vubr\$ is a direct cause of \$Uvzi\$ and is not a confounder of the \$Zuph \textbackslash{}to Uvzi\$ relationship (there is no backdoor path from \$Zuph\$ to \$Uvzi\$ through \$Vubr\$ that needs blocking), conditioning on \$Vubr\$ would block the causal influence if \$Vubr\$ were a mediator, but here \$Vubr\$ is a separate cause. However, the key issue is that \$Zuph\$ and \$Vubr\$ share a common cause \$Wibl\$. If we condition on \$Vubr\$, we induce a collider bias or simply restrict the sample in a way that does not represent the general causal mechanism unless we are specifically isolating a direct effect while holding a confounder constant. But \$Vubr\$ is not a confounder of \$Zuph\$ and \$Uvzi\$.} \\
\\ 
\texttt{Actually, let's re{-}evaluate the path: \$Zuph \textbackslash{}to Uvzi\$ is the direct effect. \$Vubr \textbackslash{}to Uvzi\$ is another direct effect. \$Wibl\$ causes both \$Zuph\$ and \$Vubr\$. There is no path from \$Zuph\$ to \$Uvzi\$ that goes through \$Vubr\$. Therefore, \$Vubr\$ is not a} \\
            }
        }
    }
    & & \\ \\

    \theutterance \stepcounter{utterance}  
    & & & \multicolumn{2}{p{0.3\linewidth}}{
        \cellcolor[rgb]{0.9,0.9,0.9}{
            \makecell[{{p{\linewidth}}}]{
                \texttt{\tiny{[GM$|$GM]}}
                \texttt{No} \\
            }
        }
    }
    & & \\ \\

    \theutterance \stepcounter{utterance}  
    & & & \multicolumn{2}{p{0.3\linewidth}}{
        \cellcolor[rgb]{0.9,0.9,0.9}{
            \makecell[{{p{\linewidth}}}]{
                \texttt{\tiny{[GM$|$GM]}}
                \texttt{yes} \\
            }
        }
    }
    & & \\ \\

    \theutterance \stepcounter{utterance}  
    & & & \multicolumn{2}{p{0.3\linewidth}}{
        \cellcolor[rgb]{0.9,0.9,0.9}{
            \makecell[{{p{\linewidth}}}]{
                \texttt{\tiny{[GM$|$GM]}}
                \texttt{game\_result = LOSE} \\
            }
        }
    }
    & & \\ \\

    \theutterance \stepcounter{utterance}  
    & & \multicolumn{4}{p{0.6\linewidth}}{
        \cellcolor[rgb]{0.95,0.95,0.95}{
            \makecell[{{p{\linewidth}}}]{
                \texttt{\tiny{[P1$\leftrightarrow$P]}}
                \texttt{Post{-}episode Reflexion (not sent to GM):} \\
\texttt{New plan: I failed because I overcomplicated the causal graph by inventing a backdoor path and misidentifying \$Vubr\$ as a confounder, leading to an incorrect "No" answer. The prompt explicitly states \$Wibl\$ affects both \$Zuph\$ and \$Vubr\$, but crucially, there is no path from \$Zuph\$ to \$Uvzi\$ that goes through \$Vubr\$. Since \$Vubr\$ is a direct cause of \$Uvzi\$ independent of \$Zuph\$ (other than sharing a common cause \$Wibl\$), conditioning on \$Vubr\$ in Method 1 does not block the direct causal path \$Zuph \textbackslash{}to Uvzi\$ nor does it open a new biasing path; it simply stratifies the effect. However, the core error was assuming Method 1 was superior due to confounding. In reality, to estimate the *total* causal effect of \$Zuph\$ on \$Uvzi\$, one should look at the general correlation (Method 2) without conditioning on a parallel cause (\$Vubr\$) that shares a common ancestor, as conditioning might induce collider bias or unnecessary restriction. The correct answer is "No" because Method 1 (conditioning on \$Vubr\$) is not more correct than Method 2 for estimating the total effect; in fact, Method 2 is the standard approach for total effect unless specific direct/indirect decomposition is requested. Next time, I will strictly draw the DAG: \$Wibl \textbackslash{}to Zuph\$, \$Wibl \textbackslash{}to Vubr\$, \$Zuph \textbackslash{}to Uvzi\$, \$Vubr \textbackslash{}to Uvzi\$. I will identify that \$Vubr\$ is a parallel cause, not a confounder of the \$Zuph \textbackslash{}to Uvzi\$ link. I will conclude that Method 2 is sufficient and correct for the total effect, while Method 1 is unnecessary and potentially biased, thus answering "No" with clear reasoning about the absence of a backdoor path requiring adjustment.} \\
            }
        }
    }
    & \\ \\

\end{supertabular}
}

\end{document}
